\documentclass[11pt]{article}
\usepackage[margin=1in]{geometry}
\usepackage{amsmath,amssymb,amsthm}
\usepackage[hidelinks]{hyperref}
\newtheorem{proposition}{Proposition}
\newtheorem{lemma}{Lemma}
\title{Two explicit bridges: centred weighted Gaussian phases and simple-quotient counts}
\author{Coordinator working deductions}
\date{16 November 2026}
\begin{document}
\maketitle
\noindent\textbf{Status.}
These are unregistered working deductions, conditional where stated on
existing internal proofs. No independent review, formal verification or
novelty is claimed. No agent goals or authoritative status records are changed.

\section{A centred weighted Gaussian corollary}
Use the exact hypotheses of
\href{../moment-vanishing-audit/SEMIPOSITIVE_WEIGHTED_GAUSSIAN_REDUCTION_AND_SAFE_NORMAL_FORMS.md}
{the nonnegative weighted Gaussian reduction}: a nonsingular Gaussian
quadratic form and a phase both have weight two in supplied nonnegative
rational linear-coordinate weights. In its notation,
\[
 P=W\cdot H(Z)+F(Z,T),\qquad g=tH(g).
\]
The residual polynomial $F$ has no constant or linear term in the
positive-weight variables $T$. Over every commutative $\mathbb Q$-algebra,
the internal theorem says that, under purity, all-marker safety is
equivalent to polynomiality of all the coordinate series $g_i$.

\begin{proposition}
Assume the cited internal reduction theorem. If in addition the ordinary
linear part of $P$ vanishes, every pure such phase is safe for all original
polynomial markers, over every coefficient $\mathbb Q$-algebra.
\end{proposition}
\begin{proof}
The ordinary linear terms in the displayed phase are precisely
$W\cdot H(0)$. Thus $H(0)=0$. The zero vector solves $g=tH(g)$, and
coefficient recursion gives uniqueness of its solution in $tR[[t]]^n$.
Hence $g=0$, and the cited safety criterion applies.
\end{proof}

In particular, an even phase has no ordinary linear part. The homogeneous
four-Gaussian invariant lift of an exceptional Haar phase is even.
Consequently the radial-transfer theorem gives pure-to-safe inclusion for
every exceptional phase whose actual homogeneous lift admits the supplied
nonnegative weighting. Equivalently, the internal minimum bad-marker
degree bounds (eight for $2O$, twelve for $2I$) contradict the linear
witness required by the weighted theorem.

\textbf{Coverage limitation.}
This does not prove that unrestricted exceptional phases admit such weights.
It does not establish that this corollary extends already closed torus
families. In the homogeneous invariant setting, linear orthogonal grading
symmetries should be compared with the existing left-torus classification
before claiming a larger class or new priority. The algebraic Gaussian
periods here are exponential generating functions; they must not be
substituted for the ordinary Cauchy resolvents in the resolvent-splitting
criterion.

\section{A sharper finite-quotient rank test}
Let $\operatorname{rk}H$ denote the minimum number of group generators.

\begin{lemma}
If a group $H$ has twenty distinct normal subgroups with quotient $A_5$,
then $\operatorname{rk}H\ge3$.
\end{lemma}
\begin{proof}
A two-generated group has at most as many epimorphisms onto $A_5$ as
$A_5$ has ordered generating pairs, namely $2280$.
Each distinct quotient kernel contributes exactly
$|\operatorname{Aut}(A_5)|=120$ epimorphisms.
Twenty kernels would therefore contribute $2400>2280$, a contradiction.
The same conclusion follows from Hall's direct-power formula:
$A_5^{19}$ is two-generated but $A_5^{20}$ is not.
\end{proof}

The classical threshold and Hall formula are recorded in
\href{https://seis.bristol.ac.uk/~tb13602/docs/BC_1.pdf}
{Burness--Crestani, On the generating graph of direct powers of a simple
group}, introduction, page two. An exact permutation enumeration of all
sixty even permutations on five letters independently gave $2280$ ordered
generating pairs, $120$ conjugation automorphisms and $19$ orbits.
These finite counts do not certify any knot representation.

\begin{proposition}
Assume the internal branched-cover bound
$\operatorname{mr}(K)\ge1+\operatorname{rk}\pi_1(\Sigma_2K)$.
If $\pi_1(\Sigma_2K)$ has twenty distinct $A_5$ quotient kernels and $K$
has a certified four-bridge presentation, then
$\operatorname{mr}(K)=b(K)=4$.
\end{proposition}
\begin{proof}
The lemma and cover bound give $\operatorname{mr}(K)\ge4$.
The four-bridge certificate gives $\operatorname{mr}(K)\le b(K)\le4$.
\end{proof}

\textbf{Application limitation.}
The twenty quotient kernels must all belong to the same double-cover group,
and are distinguished modulo target automorphisms, not just conjugacy.
The existing first $A_5$ maps on twelve different knots do not meet this
threshold. No new knot case is certified by this note. A complete count or
twenty retained distinct-kernel certificates would be the next obligation;
a broad representation search has not been launched. This refines the
existing all-homomorphism counting bound by counting epimorphisms and their
kernels instead.

\end{document}
